\chapter*{Déclaration de l'utilisation de l'intelligence artificielle générative}
\addcontentsline{toc}{chapter}{Déclaration de l'utilisation de l'intelligence artificielle générative}

\textcolor{eclipsepurple}{\textbf{Le texte en violet est informatif et doit être supprimé de la version finale du document.} L'utilisation de l'intelligence artificielle générative doit obligatoirement être déclarée. Il est également fortement suggéré d'indiquer si elle n'a pas été utilisée. L'étudiante ou l'étudiant doit obligatoirement prendre connaissance du \href{https://libguides.biblio.usherbrooke.ca/ld.php?content_id=37838948}{guide sur l’utilisation responsable de l'intelligence artificielle générative aux cycles supérieurs}, produit par Myriam Beaudet et Mireille Léger-Rousseau du Service des bibliothèques et archives de l'Université de Sherbrooke, dans lequel il est mentionné que\,:}

\begin{customquote}
\textcolor{eclipsepurple}{%
«\,Toute utilisation d'IA ayant contribué au contenu d’une production doit être déclarée. Cette transparence maintient le lien de confiance et permet de se conformer aux obligations de déclaration. La déclaration doit généralement préciser l'usage fait de l'IA et identifier l'outil utilisé ainsi que sa version. Les usages prévus doivent être discutés avec la ou les directions de recherche. […] L'usage non transparent ou non autorisé constitue un manquement à l'intégrité académique et peut entraîner des sanctions selon le Règlement des études (Règlement 2575-009)\,».%
}
\end{customquote}

\textcolor{eclipsepurple}{La déclaration doit comprendre les éléments suivants\,:} 

\begin{itemize}
    \item \textbf{Outils utilisés}\,: [Inscrire le nom de l'outil et la version. Ex\,: Claude 3.5 Sonnet (Anthropic), ChatGPT-4 (OpenAI), GitHub Copilot].
    \item \textbf{Nature de l'assistance}\,: [Indiquer les types de tâches déléguées. Ex\,: révision linguistique, génération de code Python, traduction anglais-français, génération de visualisations, reformulation de sections. Consulter au besoin \href{https://libguides.biblio.usherbrooke.ca/ld.php?content_id=37838948}{l’annexe C du guide sur l’utilisation responsable de l'intelligence artificielle générative aux cycles supérieurs}].
    \item \textbf{Étendue}\,: [Indiquer les chapitres ou sections concernés. Ex\,: Chapitre 2 - Revue de littérature, Section 3.2 - Analyse des données, figures 4.1 à 4.3.].
    \item \textbf{Type de données fournies à l'IA}\,: [Décrire les types de contenus fournis ET confirmer qu'aucune donnée confidentielle, ce qui inclut des renseignements personnels, et qu’aucune donnée non publiée ou couverte par une entente de confidentialité n'a été transmise, OU décrire les autorisations obtenues].
    \item \textbf{Contrôle et validation}\,: [Décrire les procédures de vérification\,: relecture, validation des références, tests, réplication des analyses, consultation des sources primaires, etc.].
    \item \textbf{Mesures de protection}\,: [Si pertinent, décrire les procédures d'anonymisation, l'utilisation de modèles locaux, etc.].
\end{itemize}











